\section{Introduction}
\label{sec:introduction}

A determinant can have an Euler product without its primitive factors being
rational primes. Periodic-point zeta functions organize closed dynamical
orbits, while arithmetic Euler products organize prime objects of a specified
ring or field. The two forms may look similar, but a factorwise
identification must also preserve the type of the primitive object, its
clock, its marker, its repetitions, and the operator that owns the
determinant. Periodic-point zeta functions and finite-shift determinant
identities are foundational prior art \citep{artin1965periodic,bowen1970zeta};
their existence alone supplies no rational-prime dictionary.

This distinction is concrete for the full shift over $\F_q$. A primitive
necklace of length $n$ has the finite-field degree clock $n\log q$ and carries
the marked factor
\[
  \bigl(1-z^n q^{-ns}\bigr)^{-1}.
\]
The complete source product collapses to the reciprocal of the scalar
determinant $1-zq^{1-s}$. This ledger is exact: its degree counts agree with
the classical counts of monic irreducible polynomials over $\F_q$. Recent
necklace-polynomial work reiterates that count equality
\citep{chebolu2026necklace}. Paper 42 does not claim any of these identities
as new.

The narrower question is whether the source factors can be retyped, without
alteration, as the rational-prime factors $(1-zp^{-s})^{-1}$. The target is
not merely another indexing set. A rational prime has clock $\log p$, one
primitive marker $z$, multiplicity one in the Euler product, and a separate
diagonal operator owner. Keeping those fields explicit turns a vague
resemblance into three elementary compatibility tests.

\begin{figure}[t]
  \centering
  \resizebox{\textwidth}{!}{\begin{tikzpicture}[
  font=\small,
  >=Latex,
  node distance=10mm and 13mm,
  box/.style={draw=charcoal, rounded corners=2pt, very thick, align=center,
    minimum height=12mm, text width=37mm, fill=softgray},
  source/.style={box, draw=formalblue, fill=formalblue!8},
  control/.style={box, draw=governancepurple, fill=governancepurple!8},
  target/.style={box, draw=warningamber, fill=warningamber!9},
  owns/.style={->, very thick, draw=formalblue},
  exact/.style={->, very thick, draw=governancepurple},
  blocked/.style={->, very thick, dashed, draw=warningamber}
]
\node[source] (shift) at (0,0) {full $q$-shift\\one vertex, $q$ loops};
\node[source] (necklace) at (4.8,0)
  {$\gamma\in\Prim_q$\\length $n$, marker $z^n$\\clock $n\log q$};
\node[source] (det) at (4.8,-2.35)
  {$D_q(s,z)=1-zq^{1-s}$\\source-owned determinant};

\node[control] (ff) at (9.6,1.45)
  {monic irreducible $P\in\F_q[x]$\\degree $n$, norm $q^n$\\same counts only};
\node[target] (prime) at (9.6,-1.45)
  {rational-prime atom $p$\\marker $z$, clock $\log p$\\would require
  $p=q^n$ and $n=1$};

\draw[owns] (shift) -- (necklace);
\node[font=\scriptsize,align=center,text=charcoal] at (2.4,0.85)
  {primitive cyclic classes};
\draw[owns] (necklace) -- node[left,align=right]{Euler\\regrouping} (det);
\draw[exact] (necklace) -- (ff);
\draw[blocked] (necklace) -- (prime);

\node[draw=warningamber, very thick, circle, inner sep=1.2pt,
  fill=white] at ($(necklace)!0.43!(prime)$) {$\times$};
\node[align=center, text=charcoal, below=8mm of det, text width=118mm]
  {Purple records a countwise function-field control. The dashed lane is the
  frozen rational-prime retyping test; marker and multiplicity still fail.};
\end{tikzpicture}
}
  \caption{The full shift owns a valid function-field ledger: primitive
  necklaces have length $n$, clock $n\log q$, marker $z^n$, and norm $q^n$.
  The count equality with monic irreducible polynomials is a positive control.
  A rational-prime atom instead has clock $\log p$, marker $z$, and
  multiplicity one. The dashed attempted descent is blocked when all source
  fields are required to remain fixed.}
  \label{fig:source-target-types}
\end{figure}

The comparison yields four precise contributions.

\begin{enumerate}[label=(\roman*)]
\item We prove that no total map from all primitive source necklaces to
  rational primes preserves the exact source clock. The two-letter primitive
  class $[01]$ already forces the composite image $q^2$.
\item We prove a separate factor obstruction. Equality of marker and analytic
  weight forces length one and $p=q$, where the $q$ source factors collide
  with one rational-prime factor.
\item We compare the complete marked determinants without using zeros. Their
  logarithms have different first $z$ coefficients on the common
  absolute-convergence domain.
\item We classify six declared repairs by the coordinate they abandon. The
  classification preserves the positive function-field ledger and refuses to
  transfer ownership from a changed marker, projection, or target operator.
\end{enumerate}

The conclusion is deliberately local. It does not preclude an arbitrary
relabeling, a partial factor, an induced system, a countable prime inventory,
or an infinite-memory construction. Each such proposal may define a useful
new model, but it must declare a new object and operator contract. Nor does
the paper claim that a bounded literature search proves priority. The source
identities have zero novelty, the broad mechanism has zero novelty, and the
exact typed closure is assessed at only $3/10$.

The rest of the paper separates prior ownership and chronology
(\cref{sec:prior-scope}), fixes the two factor types
(\cref{sec:source-ledger}), proves the three independent obstructions
(\cref{sec:non-descent}), and classifies repairs
(\cref{sec:repairs}). The strict Route disposition and reproducibility
boundary appear in \cref{sec:route}; exact auxiliary calculations are collected
in \cref{app:exact-details}.
